\section{Introduction}\label{sec:intro}

Coastal salt marshes lie at the intersection of biological invasion,
sea-level change, and intensive human management. \emph{Spartina
alterniflora} Loisel., a perennial cordgrass native to the Atlantic and
Gulf coasts of North America, was deliberately introduced to China in
1979 and subsequently spread along most of the country's coastline,
altering benthic habitat, marsh productivity, and soil and ecosystem
processes
\citep{wang2018soilcarbon,zheng2018productivity,wang2010benthic}. Remote-sensing studies
document rapid, spatially heterogeneous invasion at national scale from
Landsat time series \citep{liu2018rapid,mao2019rapid}, phenology-driven
mappability that varies with latitude \citep{zhang2022latitudinal}, and
expansion modes and documented dieback
\citep{yan2023monitoring}. Management interventions such as cutting
combined with waterlogging alter regrowth under specific timing
\citep{yuan2011cutting}, so monitoring is not a one-time classification
problem: it must distinguish invasion, eradication, and recurrence over
decades.

\subsection{The core problem}
Long-term monitoring of this kind requires evidence that is comparable
across four decades and across sensor generations. The observations
themselves change character over time: Landsat~5, 7, 8, and~9 supply a
30\,m record beginning in the 1980s, while Sentinel-1 C-band SAR and
Sentinel-2 optical imagery supply 10--20\,m observations from the
mid-2010s \citep{wulder2012archive,torres2012s1,drusch2012s2}. Scenes
from different sensors differ in geometry, radiometry, revisit
cadence, cloud sensitivity, and in the quality metadata attached to
each acquisition. A single Landsat scene, a multi-tile Sentinel-2
datatake, and a Sentinel-1 IW frame do not define the same spatial
event, and treating them as interchangeable pixels silently bakes
sensor-era structure into every downstream product.

Two further complications are specific to management-oriented
monitoring. First, available labels are heterogeneous and uncertain:
national mapping products, local rule-based masks, and future
field-verified references differ in class definitions, minimum mapping
units, and licensing, and none of the labels currently available to the
project has been verified in the field. Second, intertidal vegetation
is observed through tide and phenology: inundation changes the apparent
surface between acquisitions \citep{oconnell2017tmii}, and the seasonal
window that best separates \emph{S.~alterniflora} from native
vegetation shifts with latitude \citep{zhang2022latitudinal}.

\subsection{Limitations of existing routes}
Existing Spartina mapping studies establish that the species is
observable from satellites, using high-resolution
\citep{ai2017phenology,liu2017monitoring} or national Landsat/Sentinel
records \citep{liu2018rapid,mao2019rapid,yan2023monitoring}, and
SAR--optical fusion is effective for complex coastal wetlands
\citep{jamali2022swin}. These studies are typically single-epoch or
few-epoch classifications over a defined extent, however. They do not
provide (i)~a sensor-agnostic observation contract under which a
multi-tile datatake, a Landsat frame, and a SAR pass are represented as
comparable events; (ii)~explicit provenance, quality, and rights
attached to every observation and label; or (iii)~a leakage-controlled
sampling frame that remains valid as new sensors and new regions are
added. Separately, the Earth-observation representation-learning
literature has produced temporal encoders
\citep{russwurm2020selfattn,garnot2020pixelset,yuan2022sitsformer} and
general-purpose geospatial foundation models
\citep{cong2022satmae,fuller2023croma,jakubik2024prithvi,xiong2024dofa},
but model capability does not resolve what a trustworthy, traceable
multi-decade observation actually is. Cross-region and cross-sensor
generalization remains a documented weakness when validation ignores
spatial structure \citep{tuia2016da,luo2022crossspatiotemporal,ploton2020spatial}.

\subsection{Approach}
We introduce \textbf{Spartina Earth}, a provenance-aware multisensor
framework organized around a fixed hierarchy: a provisional
regional bay envelope, a fixed metric analysis cell, and an
observation event that binds source scenes, sensor-specific quality
contracts, acquisition time, tide and phenology context, label tier,
rights, and a cryptographic provenance record. Quality is evaluated
against the cell and only against the scenes that actually contribute
to it; SAR eligibility uses the measured per-acquisition footprint
rather than a representative orbit frame; labels are held in explicit
quality tiers with redistribution rights separated from internal-use
rights; and every split is grouped by spatial units so adjacent
patches never cross train and test \citep{roberts2017cv,gebru2021datasheets,mitchell2019modelcards}.

\subsection{Contributions and scope}
This paper is a methods-and-evidence paper, not a foundation-model
paper. Its present contributions are:
\begin{enumerate}[leftmargin=*]
\item a provenance-aware observation model for cross-generation
coastal monitoring---envelope, analysis cell, acquisition group,
modality quality contracts, label tiers and rights---implemented and
tested in an operational data factory on Google Earth Engine
\citep{gorelick2017gee};
\item a three-bay (Hangzhou, Sanmen, Yueqing) metadata census of
18{,}435 physical scenes from six sensors spanning 1985--2026, with
explicit distinctions among not-operational, no-scene, and no-quality
periods;
\item three independently executed audits that correct defects in the
initial frame: false-zero label statistics from invalid geometry and
cell-edge label bleed; a cloud-gating bug that had rejected every
2022 autumn Sentinel-2 event in all three bays (31{,}135 wrongly
rejected 10\,km pairs across the full archive); and a quantified
demonstration that representative SAR footprints cannot be used for
production eligibility;
\item a controlled eight-cell, 13-product real-pixel multimodal pilot
that closes the observation chain on real bytes---server-asserted
same-datatake optical merges, the first eligible native-30\,m
Landsat~8/9 scaling validation, and an identity-decibel Sentinel-1
product independently recomputed server-side---and that establishes
nominal acquisition frames (representative SAR polygons \emph{and}
nominal MGRS tile polygons) as planning prefilters: they miss five of
eight pilot cells that real datatake footprints cover;
\item preliminary, explicitly limited Pilot-0 baselines showing both
that the data support learned models and where current evidence
collapses (tiny test support, SILVER/WEAK disagreement, seed-level
model failures); and
\item a target-independent design for national scaling that uses the
2015 national Spartina map as a stratification prior rather than as the
study boundary.
\end{enumerate}
A sensor-agnostic temporal representation model (SpartinaFM) is the
flagship \emph{future} direction; no representation-learning results
are claimed. GOLD field/UAV validation does not yet exist, bay
envelopes are provisional, and tide modeling is selected but not
deployed; these limits are carried explicitly through every section.
